\chapter{Resultados Parciais}

Até a presente etapa da pesquisa, os esforços foram concentrados na fundamentação teórica e na análise das condições necessárias para estender as decomposições de erro e diversidade de ensembles ao domínio de séries temporais. Como primeiro resultado parcial, por meio do mapeamento de abordagens sobre diversidade, evidenciou-se que a transposição direta de métricas de diversidade desenvolvidas para dados independentes e identicamente distribuídos (i.i.d) não captura os efeitos de dependência temporal, autocorrelação e da não estacionariedade.

Para além da análise teórica, conduziu-se um experimento preliminar para testar a aplicação direta da descorrelação de erros ($1 - \rho$) na seleção estática de modelos para séries temporais. O protocolo partiu do melhor preditor individual na validação, adicionando iterativamente os modelos candidatos que apresentavam a maior diversidade (menor correlação de erros) em relação à previsão agregada do conjunto já selecionado ($x_1 + x_2$, $x_1 + x_2 + x_3$, até $x_n$), combinados por média e mediana. No conjunto de testes, confirmou-se que o acréscimo de modelos guiado apenas pela maximização dessa diversidade não reduziu o erro do \textit{ensemble} em comparação ao melhor modelo isolado. Esse comportamento evidencia na prática que ignorar a relação entre o erro individual e diversidade pode comprometer o desempenho da combinação. \textcolor{red}{Consegue }

Como resultados esperados para as próximas etapas, pretende-se: 1) concluir a dedução formal para a decomposição de viés, variância e diversidade sob premissas não-i.i.d; 2) avaliar empiricamente, em conjuntos de séries temporais com características diferentes, a relação entre a diversidade mensurada e o erro de previsão ao longo do tempo; e 3) propor diretrizes quantitativas que auxiliem a seleção dinâmmica em cenários de séries temporais.